\section{The Poet and the Saint}

\begin{quotex}
The \emph{gift} of the mystic is akin to the gift of poetry.
\flright{\textsc{Caroline Stephen}, \emph{Light Arising}}

Il n'y a qu'une tristesse, c'est de n'être pas des saints.
\flright{\textsc{Léon Bloy}}

The characteristic feature of the great saints is that they get at the very heart of the matter. The most obvious things are invisible because they are concealed in human beings; nothing is harder to evince than what is self-evident. Once it is uncovered or rediscovered, it develops explosive strength. \textbf{Saint Anthony} recognized the power of the solitary man. \textbf{Saint Francis} that of the poor man. \textbf{Stirner} that of the ``only'' man. At bottom, everyone is solitary, poor, and ``unique'' in the world.
\flright{\textsc{Ernst Junger}, \emph{Eumeswil}}
\end{quotex}


\paragraph{Saint Francis}


In his \emph{Preparation of the Soul for Contemplation}, \textbf{Richard of St. Victor} reminds us that many know about contemplation, some by science, others by experience. \textbf{St. Francis} is the exemplar of the latter; Dante's biographer, Leonardo Bruni, describes him this way:

\begin{quotex}
The blessed Francis not through science, not through discipline of the schools, but by mental possession and ecstasy, applied his mind so strongly to God that he was transfigured beyond the measure of human sense, and knew more of God than the theologians know through their study or through letters.
\end{quotex}


Embracing Poverty, St. Francis's life was itself a Poem. With the Stigmata, God's presence was visceral to him. The Franciscan friars were always ``minor'', so there was no possibility to advance in grades as in an initiatic organization. God decides who belongs to Him, not some club.

Nevertheless, there were two who arose together: \textbf{Saint Francis} and \textbf{Saint Dominic}, the Poet and the Philosopher. Both are necessarily joined as explained by the Franciscan \textbf{Saint Bonaventura}. What can the philosopher learn from Saint Francis? Well, besides material poverty there is spiritual poverty. The philosopher, who is poor in spirit, must needs shed the many useless opinions that have been implanted in his soul.

There are three degrees of knowledge:

\begin{enumerate}


\item Opinion

\item Faith

\item Understanding
\end{enumerate}


Faith will spill the space previously held by opinion. Faith cannot be proved through books, since it is an act of the will. It is achieved through a spiritual path of prayer, meditation, worship, and purification. Only afterwards will understanding follow, the only kind of knowledge worthy of the philosopher.

\paragraph{Saint Anthony}


\begin{quotex}
Godliness is an abomination to the sinner.
\flright{\textsc{Saint Anthony}}

A man who has not passed through the inferno of his passions has never overcome them. They then dwell in the house next door, and at any moment a flame may dart out and set fire to his own house. Whenever we give up, leave behind, and forget too much, there is always the danger that the things we have neglected will return with added force.
\flright{\textsc{Carl Jung}}
\end{quotex}


As with Saint Francis, we learn from Saint Anthony that true understanding comes from living, not from abstract intellectual knowledge. Again like Saint Francis, Saint Anthony came from a prosperous family and could have had a cushy life had he wanted it. On the other hand, Saint Anthony chose the life of a solitary hermit rather than life in a brotherhood.

Yet that does not make the spiritual life easier. The devil fought Anthony by afflicting him with \emph{acedia}, boredom, laziness, and \emph{fantasies of women}. If you don't recognize those temptations, or worse, think that they are insignificant, then your spiritual life will be stunted. Whoever is not so tempted cannot be saved. St Anthony said:

\begin{quotex}
This is the great work of a man: always to take the blame for his own sins before God and to expect temptation to his last breath. Whoever has not experienced temptation cannot enter into the Kingdom of Heaven. Without temptations no one can be saved.
\end{quotex}


During a self-imposed lockdown, there is the opportunity to experience the hermit life. Especially for those who are most vulnerable, the awareness of the possibility of death by an invisible bug is palpable. That could be to our advantage. Saint Anthony elaborates:

\begin{quotex}
If we so live as people dying daily, we will not commit sin. As we rise daily, let us suppose that we shall not survive until evening and as we prepare for sleep, let us consider that we shall not awaken. Our life is uncertain. If we live this way, we will not sin, nor will we crave anything, nor bear a grudge against anyone.
\end{quotex}


Many ``seekers'' search for the next book, or idea, or secret society, but that is a form of evasion. Saint Anthony reminds us that, ``The Kingdom of Heaven is within you.'' That refers to our intellectual soul, which needs to be rightly ordered. In his words:

\begin{quotex}
All virtue needs is our willing, since it is in us and arises from us. For virtue exists when the soul maintains its intellectual part according to nature.
\end{quotex}


Finally, we need to avoid the temptation of dialectics. Saint Anthony explained this to some seekers:

\begin{quotex}
You want us not to worship God without demonstrated proof from arguments? Does knowledge of God come through arguments or an act of faith? ... Faith comes from the disposition of the soul, but dialectic is from the skill of those who construct it.
\end{quotex}


Once again we see that understanding follows faith. Do not be deceived by the skill of the sophists.

\paragraph{Saint Max}

\begin{quotex}
As this rose is a true rose to begin with, this nightingale always a true nightingale, so I am not for the first time a true man when I fulfill my calling, live up to my destiny, but I am a ``true man'' from the start. My first babble is the token of the life of a ``true man'', the struggles of my life are the outpourings of his force, my last breath is the last exhalation of the force of the ``man''. The true man does not lie in the future, an object of longing, but lies, existent and real, in the present. 
\flright{\textsc{Max Stirner}, \emph{The Ego and its Own}}
\end{quotex}

Karl Marx in \emph{The Holy Family} and \emph{German Ideology} ironically referred to Max Stirner as Saint Max. That was because Marx saw him as an obstacle to his revolution. For Marx, not to mention certain occult teachings based upon some law of evolution, the ``true man'' will appear in some distant future.

Ownership implies the power to do what one wills with one's property. So self-ownership, one's ``own'', implies the same. But it is more than just ownership over one's body, it requires self-mastery over desires, emotions, and thoughts. Junger offers this suggestion:

\begin{quotex}
If one manages to separate essence from flesh, if one manages, that is, to get distance from oneself, then one climbs the first step toward spiritual power. Many exercises are geared to this---from the soldier's drill to the hermit's meditation.
\flright{Venator in \textit{Eumeswil}}
\end{quotex}


We could add that today perhaps extreme sports, or even bodybuilding which seems to be a sudden craze in some circles, are means to overcome the lethargy of the body. Yet the physical is just one phase, the childhood of the spiritual life. Saint Max describes the phases:

\begin{itemize}


\item \textbf{Childhood: Realism}. Encounter with the world as other. The world as material.

\begin{quotex}
the child from the outset lives a life of struggle against the entire world, it resists everything and everything resists it.
\end{quotex}


\item \textbf{Youth: Idealism}. Awareness of spirit. The task of youth is

\begin{quotex}
to bring to light \emph{pure thought}, to devote oneself to it --- in this is the \emph{joy of youth}, and all the bright images of the world of thought --- truth, freedom, mankind, \emph{Man}, etc. --- illumine and inspire the youthful soul.
\end{quotex}


\item \textbf{Adult: Egoism}. The full development, and hence responsibility, of the I. Saint Max elucidates:

\begin{quotex}
How I find myself behind the \emph{things}, and indeed as \emph{spirit}, so subsequently, too, I must find myself behind the \emph{thoughts}, i.e., as their creator and owner. In the period of spirits, thoughts outgrew me although they were the offspring of my brain; like delirious fantasies they floated around me and agitated me greatly, a dreadful power. The thoughts became themselves \emph{corporeal}, they were spectres like God, the Emperor, the Pope, the Fatherland, etc.; by destroying their corporeality, I take them back into my own corporeality and \emph{announce}: I alone am corporeal. And now I take the world as it is for me, as my world, as my property: I relate everything to myself.
\end{quotex}

\end{itemize}


Saint Max manages a level of detachment from thoughts, experiencing them as mere things in his mind. Otherwise, those spectres become reified and one sees only the sclerotic remains of those ideas. The saint must take responsibility and ownership for his worldview and way of life. Then one becomes a spirit in the world. Saint Max makes it clear:

\begin{quotex}
If I as spirit rejected the world with the deepest contempt for it, then I as proprietor reject the spectres or ideas into their emptiness. They no longer have power over me, just as no `earthly force' has power over the spirit.
\end{quotex}


Yet Max Stirner stops short. Having expunged the false idols of God in his mind, he remains satisfied with his own Ego. That should have been the impulse to seek the True God, even if only to recognize that Atman (ego) is Brahman (God). Ultimately, Max did not own his own birth nor his own death.


\flrightit{Posted on 2020-08-09 by Cologero}

